\documentclass[11pt,a4paper]{article}

\usepackage[utf8]{inputenc}
\usepackage[margin=2.4cm]{geometry}
\usepackage{amsmath,amssymb}
\usepackage{booktabs}
\usepackage{array}
\usepackage{enumitem}
\usepackage{listings}
\usepackage[hidelinks]{hyperref}

\setlist[itemize]{label=$\bullet$}
\setlength{\parindent}{0pt}
\setlength{\parskip}{0.55em}

\lstset{
  language=Java,
  basicstyle=\ttfamily\footnotesize,
  keywordstyle=\bfseries,
  commentstyle=\itshape,
  columns=fixed,
  basewidth=0.525em,
  keepspaces=true,
  upquote=false,
  literate={'}{{\char13}}1,
  showstringspaces=false,
  frame=single,
  framerule=0.4pt,
  xleftmargin=0.5em,
  xrightmargin=0.5em,
  aboveskip=0.6em,
  belowskip=0.6em,
  breaklines=true
}
\lstdefinestyle{out}{
  language={},
  keywordstyle={},
  commentstyle={},
  frame=leftline,
  framerule=0.8pt,
  xleftmargin=0.6em,
  xrightmargin=0pt
}
\newcolumntype{S}[1]{>{\ttfamily\small\raggedright\arraybackslash}p{#1}}
\newcolumntype{D}[1]{>{\raggedright\arraybackslash}p{#1}}
\newcommand{\code}[1]{\lstinline[basicstyle=\ttfamily\normalsize,basewidth=0.5em]{#1}}

% Questions: Q<exercise>.<n>
\newlist{questions}{enumerate}{1}
\setlist[questions]{label=\textbf{Q\theexnum.\arabic*}, leftmargin=3.2em, itemsep=0.5em}
\newcounter{exnum}

\newcommand{\exercise}[3]{%
  \clearpage
  \stepcounter{exnum}%
  \section*{Exercise \theexnum\quad #1}
  \addcontentsline{toc}{section}{Exercise \theexnum: #1}
  \textit{Difficulty: #2.\quad Indicative time: #3.}\par\medskip}

\newcommand{\nl}{\texttt{\char13\char92 n\char13}}
\newcommand{\spec}[1]{\par\smallskip\textbf{#1}\par}

\title{Introduction to Object-Oriented Programming (Java)\\[0.3em]
\large Practice set for weeks 1--2 (Lessons 1--9)}
\author{}
\date{}

\begin{document}
\maketitle
\vspace{-3em}

\section*{How to use this sheet}

There are three exercises of increasing difficulty. Each one consists of
\begin{enumerate}[label=(\alph*), itemsep=0pt]
  \item a \emph{specification} to implement;
  \item one or more \emph{test programs} with their expected output. Where the
        output is given, it must match yours character for character;
  \item \emph{written questions} (Q). Answer each in a few sentences, and
        justify every claim. ``It works when I run it'' is not a justification.
        Answer the questions marked \emph{before coding} before writing any code.
\end{enumerate}

\paragraph{Allowed toolbox.} Everything from \emph{Initiation à la
programmation en Java} (BOOC~I) and Lessons 1--9 of the OOP course. The
following come later in the course and are \textbf{not} allowed: inheritance
(\code{extends}), interfaces, exceptions (\code{throw}, \code{try}), and
\code{static} attributes other than constants (\code{private static final}).
Static methods are only allowed in the class that contains \code{main}, as
in BOOC~I. Every attribute is \code{private}.

\paragraph{Conventions.}
\begin{itemize}[itemsep=0pt]
  \item \code{equals} has the signature seen in Lesson~9:
        \code{public boolean equals(C other)}, where \code{C} is the class itself. It returns
        \code{false} when \code{other} is \code{null}.
  \item Because \code{equals(Object)} has not been covered yet, the \code{ArrayList}
        methods \code{contains}, \code{indexOf} and \code{remove(Object)} will
        \emph{not} call your \code{equals}. Write the loops yourself.
  \item Exceptions are not available yet, so invalid operations report failure
        through their return value (\code{false}, \code{null} or \code{-1}),
        exactly as specified.
\end{itemize}

\paragraph{What is tested.} {\small $\bullet$ = central, $\circ$ = used.}

\begin{center}
\begin{tabular}{l c c c c}
\toprule
Concept & Lesson & Ex.\,1 & Ex.\,2 & Ex.\,3 \\
\midrule
Classes, attributes, methods, \code{this}, masking   & 2--3 & $\circ$ & $\circ$ & $\circ$ \\
\code{public}/\code{private}, getters and setters    & 3    & $\bullet$ & $\circ$ & $\circ$ \\
Encapsulation: class invariants, interface design    & 4--5 & $\bullet$ & $\bullet$ & $\bullet$ \\
Constructors, overloading, chaining with \code{this(...)} & 6--7 & $\bullet$ & $\bullet$ & $\circ$ \\
Default constructor rules                            & 7    & $\bullet$ & $\circ$ & \\
Copy constructor; shallow vs.\ deep copy             & 8    & $\circ$ & $\bullet$ & $\bullet$ \\
References, aliasing, \code{null}, \code{==}         & 3, 9 & $\circ$ & $\bullet$ & $\bullet$ \\
\code{toString}, \code{equals}                       & 9    & $\bullet$ & $\bullet$ & $\bullet$ \\
End of life, garbage collection                      & 9    &         & $\circ$ & $\circ$ \\
Arrays, \code{ArrayList}, \code{String} (BOOC I)     & --   & $\circ$ & $\bullet$ & $\bullet$ \\
\bottomrule
\end{tabular}
\end{center}

%======================================================================
\exercise{Exact fractions and the hunt for $\pi$}{easy}{2 hours}

In Java, \code{0.1 + 0.2 == 0.3} is \code{false}, because a \code{double}
cannot store most decimal fractions exactly. In this exercise you build an
exact \code{Fraction} type, then use it to compute harmonic numbers exactly
and to find good rational approximations of $\pi$.

\subsection*{1.1\quad The class \texttt{Fraction}}

A \code{Fraction} stores its numerator and denominator in two
\code{private final long} attributes. Every \code{Fraction} object must
satisfy the following \emph{class invariant} from the moment its constructor
returns:
\begin{enumerate}[label=(I\arabic*), itemsep=0pt]
  \item the denominator is strictly positive;
  \item numerator and denominator have no common divisor other than $1$
        (so zero is stored as $0/1$).
\end{enumerate}
\code{Fraction} objects are \emph{immutable}: no method ever changes the
attributes after construction, and there are no setters.

\spec{Constructors}
\begin{itemize}[itemsep=2pt]
  \item \code{Fraction(long n, long d)} builds $n/d$ in normalised form. If
        \code{d == 0}, it prints exactly
        \code{Error: zero denominator, fraction set to 0} on its own line and builds $0/1$.
  \item \code{Fraction(long n)} builds the integer $n$. \emph{Must delegate} with \code{this(...)}.
  \item \code{Fraction()} builds $0$. \emph{Must delegate} with \code{this(...)}.
  \item \code{Fraction(Fraction other)} is the copy constructor.
\end{itemize}

\spec{Methods}

\begin{tabular}{@{}S{0.40\textwidth} D{0.55\textwidth}@{}}
long getNumerator(), long getDenominator() & accessors. \\
Fraction add(Fraction o) & $\mathit{this} + o$ as a \emph{new} object. Never modifies \code{this} or \code{o}. \\
Fraction subtract(Fraction o) & $\mathit{this} - o$. \\
Fraction multiply(Fraction o) & $\mathit{this} \times o$. \\
Fraction divide(Fraction o) & $\mathit{this} / o$; returns \code{null} if $o$ is zero. \\
Fraction reciprocal() & $1/\mathit{this}$; returns \code{null} if \code{this} is zero. \\
boolean isLessThan(Fraction o) & exact comparison $\mathit{this} < o$. No \code{double} allowed. \\
double toDouble() & approximate value. \\
String toString() & \code{"n/d"}, or just \code{"n"} when the denominator is $1$. \\
boolean equals(Fraction o) & same rational number. \\
\end{tabular}

Write the greatest-common-divisor computation as a \code{private} helper
method (Euclid's algorithm). You may assume that no test overflows \code{long}.

\spec{Test 1}
\begin{lstlisting}
class FractionTest {
    public static void main(String[] args) {
        Fraction a = new Fraction(6, -8);
        Fraction b = new Fraction(1, 4);
        System.out.println(a + " " + b + " " + new Fraction(10, 5)
                           + " " + new Fraction(0, -7) + " " + new Fraction());
        System.out.println(a.add(b) + " " + a.subtract(b) + " "
                           + a.multiply(b) + " " + a.divide(b));
        System.out.println(a.divide(new Fraction()));
        System.out.println(new Fraction(2, 4).equals(new Fraction(-3, -6)));
        System.out.println(a.equals(null) + " " + a.isLessThan(b) + " " + b.isLessThan(a));
        Fraction c = new Fraction(a);
        System.out.println((c == a) + " " + c.equals(a));
        System.out.println(new Fraction(3, 0));
    }
}
\end{lstlisting}
Expected output:
\begin{lstlisting}[style=out]
-3/4 1/4 2 0 0
-1/2 -1 -3/16 -3
null
true
false true false
false true
Error: zero denominator, fraction set to 0
0
\end{lstlisting}

\subsection*{1.2\quad Harmonic numbers}

In a class \code{FractionLab} (containing \code{main}), compute and print the
harmonic numbers $H_n = 1 + \frac12 + \dots + \frac1n$ exactly for
$n = 1,\dots,10$, followed by $H_{20}$:
\begin{lstlisting}[style=out]
H_1 = 1
H_2 = 3/2
H_3 = 11/6
H_4 = 25/12
H_5 = 137/60
H_6 = 49/20
H_7 = 363/140
H_8 = 761/280
H_9 = 7129/2520
H_10 = 7381/2520
H_20 = 55835135/15519504
\end{lstlisting}

\subsection*{1.3\quad Continued fractions}

A (simple) continued fraction is written $[a_0; a_1, \dots, a_k]$ and means
\[
  [a_0; a_1, \dots, a_k] \;=\; a_0 + \cfrac{1}{a_1 + \cfrac{1}{\ddots + \cfrac{1}{a_k}}}
  \qquad (a_0 \in \mathbb{Z},\; a_i \geq 1 \text{ for } i \geq 1).
\]
For example $[4; 2, 6, 7] = 415/93$.

\begin{enumerate}[label=(\alph*)]
  \item In \code{FractionLab}, write the static function
        \code{static Fraction fromContinuedFraction(long[] a)}, which evaluates
        $[a_0; a_1, \dots, a_k]$ exactly. You may assume that \code{a.length >= 1}
        and that \code{a[i] >= 1} for \code{i >= 1}.
  \item In \code{Fraction}, write the method
        \code{ArrayList<Long> toContinuedFraction()}, which returns the
        \emph{canonical} expansion of \code{this}. Let $a_0 = \lfloor x \rfloor$. If
        $x$ is an integer, the expansion is $[a_0]$. Otherwise it is $a_0$
        followed by the expansion of $1/(x - a_0)$. Careful: Java's integer
        division truncates toward zero, which is not $\lfloor\cdot\rfloor$ for negative
        numbers. You may use \code{Math.floorDiv(long, long)}.
\end{enumerate}

\spec{Test 2}
\begin{lstlisting}
System.out.println(FractionLab.fromContinuedFraction(new long[]{4, 2, 6, 7}));
System.out.println(FractionLab.fromContinuedFraction(new long[]{3, 7, 15, 1}));
System.out.println(FractionLab.fromContinuedFraction(new long[]{-3, 1, 2}));
System.out.println(FractionLab.fromContinuedFraction(new long[]{5}));
System.out.println(new Fraction(415, 93).toContinuedFraction());
System.out.println(new Fraction(355, 113).toContinuedFraction());
System.out.println(new Fraction(-7, 3).toContinuedFraction());
System.out.println(new Fraction(5).toContinuedFraction());
System.out.println(new Fraction(1, 3).toContinuedFraction());
System.out.println(new Fraction(0).toContinuedFraction());
\end{lstlisting}
\begin{minipage}[t]{0.45\textwidth}
Expected output:
\begin{lstlisting}[style=out]
415/93
355/113
-7/3
5
[4, 2, 6, 7]
[3, 7, 16]
[-3, 1, 2]
[5]
[0, 3]
[0]
\end{lstlisting}
\end{minipage}

\begin{enumerate}[label=(\alph*), resume]
  \item \textbf{Convergents of $\pi$.} We have
        $\pi = [3; 7, 15, 1, 292, 1, 1, \dots]$. The $k$-th \emph{convergent}
        is the fraction given by the first $k$ terms. Print the first seven
        convergents, each with its decimal value and its absolute error
        $|c - \pi|$ (use \code{Math.PI}). The layout is up to you, but the fractions
        must be the following:
\begin{lstlisting}[style=out]
3              3.0000000000  error 1.42e-01
22/7           3.1428571429  error 1.26e-03
333/106        3.1415094340  error 8.32e-05
355/113        3.1415929204  error 2.67e-07
103993/33102   3.1415926530  error 5.78e-10
104348/33215   3.1415926539  error 3.32e-10
208341/66317   3.1415926535  error 1.22e-10
\end{lstlisting}
  \item \textbf{Record hunt.} For every denominator $d = 1, \dots, 1000$, let
        \code{n = Math.round(Math.PI * d)} and consider $n/d$. Print each
        fraction whose error is \emph{strictly} smaller than the error of every
        fraction seen before. You should find exactly these 14 records:
\begin{lstlisting}[style=out]
3  13/4  16/5  19/6  22/7  179/57  201/64  223/71  245/78  267/85
289/92  311/99  333/106  355/113
\end{lstlisting}
\end{enumerate}

\subsection*{1.4\quad Questions}
\begin{questions}
  \item \emph{(Lessons 6--7.)} List, in execution order, every constructor
        body that runs during \code{new Fraction()}. Next, suppose you deleted every
        constructor except \code{Fraction(long, long)}. Which expressions in Test~1
        would no longer compile? Why does Java not supply a default constructor
        in that case?
  \item \emph{(Invariant.)} Your \code{equals} compares attributes one by one.
        Explain why this is correct \emph{only because of} (I1) and (I2). Give
        two attribute pairs representing $\tfrac12$ that a non-normalising
        constructor could produce, and that \code{equals} would wrongly call
        different. Then give a pair violating (I1) for which your
        \code{isLessThan} gives a wrong answer, and show the computation.
  \item \emph{(Encapsulation and aliasing.)} A classmate removes \code{final} and adds
        \begin{center}\texttt{public void setNumerator(long n) \char123\ num = n; \char125}\end{center}
        (a)~Write a client snippet of two or three lines after which
        \code{equals} returns a wrong result.
        (b)~With this setter, what does \code{Fraction b = a; a.setNumerator(5);}
        do to \code{b}? Draw the memory diagram.
        (c)~Why is \code{Fraction b = a;} harmless in \emph{your} class? Given
        that, is your copy constructor useful at all?
  \item \emph{(\code{null}.)} What happens when
        \code{a.divide(new Fraction()).add(b)} runs, and on which exact
        sub-expression? Compare the strategy of \code{divide} (return \code{null})
        with that of the constructor (print a message and build $0/1$). Which
        failure is easier to notice and debug, and why? Neither is fully
        satisfactory. Which later chapter of the course addresses this?
  \item \emph{(\code{==} vs.\ \code{equals}.)} Predict and justify the values of
        \code{new Fraction(1, 2) == new Fraction(2, 4)},
        \code{new Fraction(1, 2).equals(new Fraction(2, 4))},
        and \code{(c == a)} in Test~1.
  \item \emph{(Continued fractions.)} \code{fromContinuedFraction} returns
        $355/113$ for both $\{3, 7, 15, 1\}$ and $\{3, 7, 16\}$. Explain why.
        Does this mean \code{toContinuedFraction} is wrong? Also, what would your method
        return for $-7/3$ if it used \code{/} instead of the floor? Derive it by hand.
  \item \emph{(Optional, mathematical.)} Only four of the 14 records are
        convergents. Look at the records from $22/7$ to $333/106$ and compute
        the successive differences of their numerators and denominators.
        Explain the pattern using the \emph{mediant}
        $\frac{a}{b} \oplus \frac{c}{d} = \frac{a+c}{b+d}$, and relate the number of
        intermediate records to the term $15$ in the expansion of $\pi$.
\end{questions}

%======================================================================
\exercise{Turtle graphics on an ASCII canvas}{medium}{3--4 hours}

A \emph{turtle} is a small robot holding a pen. It walks over a sheet of paper
and leaves a trace wherever it goes while the pen is down. You will build the
sheet (\code{Canvas}), the robot (\code{Turtle}), and an interpreter for a
small drawing language, which you will then use to draw spirals and fireworks
in the terminal.

\paragraph{Coordinates.} Cell $(x, y)$ is in column $x$ (from $0$ on the left)
and row $y$ (from $0$ at the \emph{top}). ``North'' means \emph{up}, that is,
decreasing $y$. Headings are multiples of $45^\circ$, measured clockwise from
North:

\begin{center}
\small\setlength{\tabcolsep}{4pt}
\begin{tabular}{l *{8}{c}}
\toprule
heading & 0 & 45 & 90 & 135 & 180 & 225 & 270 & 315 \\
direction & N & NE & E & SE & S & SW & W & NW \\
$(\Delta x, \Delta y)$ & $(0,-1)$ & $(1,-1)$ & $(1,0)$ & $(1,1)$ & $(0,1)$ & $(-1,1)$ & $(-1,0)$ & $(-1,-1)$ \\
\bottomrule
\end{tabular}
\end{center}

\subsection*{2.1\quad The class \texttt{Canvas}}

A canvas is a rectangular grid of characters. Empty cells contain the constant
\code{BLANK = '.'} (a \code{private static final char}). The grid is stored in
a private two-dimensional \code{char} array that is \emph{never} exposed
outside the class.

\begin{tabular}{@{}S{0.40\textwidth} D{0.55\textwidth}@{}}
Canvas(int width, int height) & all cells blank. Any dimension smaller than $1$ is replaced by $1$. \\
Canvas() & a $21 \times 11$ canvas. Must delegate with \code{this(...)}. \\
Canvas(Canvas other) & an independent copy. \\
int getWidth(), int getHeight() & dimensions. \\
boolean isInside(int x, int y) & whether $(x,y)$ is a cell of the canvas. \\
boolean plot(int x, int y, char c) & writes \code{c} at $(x,y)$. Returns \code{false} and does nothing if $(x,y)$ is outside. \\
char charAt(int x, int y) & content of $(x,y)$, or \code{BLANK} if outside. \\
void clear() & blanks every cell. \\
int inkCount() & number of non-blank cells. \\
String toString() & the rows from top to bottom, separated by \nl, with no trailing newline. \\
boolean equals(Canvas o) & same dimensions and same characters. \\
\end{tabular}

\subsection*{2.2\quad The class \texttt{Turtle}}

A turtle knows the canvas it draws on (a reference), its position, its
heading, whether its pen is down, and the character its pen writes.

\begin{tabular}{@{}S{0.40\textwidth} D{0.55\textwidth}@{}}
Turtle(Canvas canvas, int x, int y) & heading $0$, pen down, pen character \code{'*'}. Draws nothing. \\
Turtle(Canvas canvas) & starts at the centre \code{(width/2, height/2)} (integer division). Must delegate with \code{this(...)}. Must not crash when \code{canvas} is \code{null}; it then starts at $(0,0)$. \\
Turtle(Turtle other) & same state, and \emph{the same canvas} (shared, not copied). \\
getX(), getY(), getHeading() & position, and heading in degrees in $[0, 360)$. \\
isPenDown(), penUp(), penDown() & pen state. \\
void setPen(char c) & pen character. \\
\end{tabular}

\spec{\texttt{boolean forward(int n)}}
If \code{n < 0}, return \code{false} and do nothing. Otherwise, if the pen is
down, first plot the current cell. Then repeat \code{n} times: move one cell
in the direction of the heading, and plot the new cell if the pen is down.
Return \code{true}. The turtle may leave the canvas: its position is still
updated, and out-of-canvas cells are simply not drawn. With a \code{null}
canvas the turtle moves but draws nothing. (Note that \code{forward(0)} draws
a single dot.)

\spec{\texttt{boolean turnRight(int degrees)}, \texttt{boolean turnLeft(int degrees)}}
\code{degrees} may be any multiple of $45$, including negative values and values
$\geq 360$. Otherwise, return \code{false} and do nothing.

\spec{\texttt{toString()} and \texttt{equals(Turtle o)}}
\code{toString} returns, for example, \code{Turtle at (8,1) heading 90 pen down '*'}
(or \code{pen up}). Two turtles are \code{equals} when they have the same position,
heading, pen state and pen character, \emph{and} draw on the same canvas object.

\subsection*{2.3\quad A tiny drawing language}

Add the method \code{boolean execute(String program)} to \code{Turtle}.
A program is a sequence of tokens separated by one or more spaces:

\begin{tabular}{>{\ttfamily}l l}
F n & \code{forward(n)} \\
L d \quad R d & \code{turnLeft(d)}, \code{turnRight(d)} \\
U \quad D & pen up, pen down \\
P c & \code{setPen(c)}, where \code{c} is a single character \\
REPEAT k [ \dots{} ] & execute the body $k \geq 0$ times. The brackets are separate tokens. \\
\end{tabular}

In this basic version, a \code{REPEAT} body never contains another \code{REPEAT}.
Commands run from left to right. If they all succeed, \code{execute} returns
\code{true}. At the \emph{first} erroneous command it stops and returns
\code{false}; commands already executed remain executed. A command is erroneous
if it is unknown, if its argument is missing, if \code{forward} or a turn
returns \code{false}, if the argument of \code{P} is not a single character, or
if a \code{REPEAT} has a negative count, no \code{[}, or no matching \code{]}.
The empty program and a blank program both return \code{true}. You may assume
that numeric arguments are well-formed integers (\code{Integer.parseInt}).
To split the program into tokens, you may use
\code{String[] tokens = program.trim().split(" +");}.

\spec{Test 1}
\begin{lstlisting}
class TurtleTest {
    public static void main(String[] args) {
        Canvas c = new Canvas(11, 7);
        Turtle t = new Turtle(c);
        System.out.println(t);
        t.forward(2);
        t.turnRight(90);
        t.forward(3);
        System.out.println(t);
        System.out.println(c);
        System.out.println("ink: " + c.inkCount());

        Canvas sq = new Canvas(9, 6);
        Turtle s = new Turtle(sq, 1, 4);
        System.out.println(s.execute("REPEAT 4 [ F 3 R 90 ]"));
        System.out.println(sq);
        System.out.println(s);

        Canvas dia = new Canvas(9, 9);
        Turtle d = new Turtle(dia, 4, 0);
        d.setPen('#');
        System.out.println(d.execute("R 135   REPEAT 4 [ F 4 R 90 ]"));
        System.out.println(dia);

        Canvas small = new Canvas(5, 5);
        Turtle o = new Turtle(small, 2, 2);
        System.out.println(o.execute("F 4 R 180 F 6"));
        System.out.println(small);
        System.out.println(o);
        System.out.println(o.execute("F 2 X 3") + " " + o);
        System.out.println(o.execute("R 30") + " " + o.execute("F -2") + " "
                           + o.execute("REPEAT 3 [ F 1") + " " + o.execute(""));
        System.out.println(o.execute("L 450") + " " + o);
    }
}
\end{lstlisting}
Expected output (shown in two columns: read the left column first):

\begin{minipage}[t]{0.46\textwidth}
\begin{lstlisting}[style=out]
Turtle at (5,3) heading 0 pen down '*'
Turtle at (8,1) heading 90 pen down '*'
...........
.....****..
.....*.....
.....*.....
...........
...........
...........
ink: 6
true
.........
.****....
.*..*....
.*..*....
.****....
.........
Turtle at (1,4) heading 0 pen down '*'
\end{lstlisting}
\end{minipage}\hfill
\begin{minipage}[t]{0.53\textwidth}
\begin{lstlisting}[style=out]
true
....#....
...#.#...
..#...#..
.#.....#.
#.......#
.#.....#.
..#...#..
...#.#...
....#....
true
..*..
..*..
..*..
..*..
..*..
Turtle at (2,4) heading 180 pen down '*'
false Turtle at (2,6) heading 180 pen down '*'
false false false true
true Turtle at (2,6) heading 90 pen down '*'
\end{lstlisting}
\end{minipage}

\spec{Test 2: a spiral driven from Java, and two turtles sharing one canvas}
\begin{lstlisting}
Canvas sp = new Canvas(13, 13);
Turtle p = new Turtle(sp);
for (int i = 1; i <= 12; ++i) {
    p.forward(i);
    p.turnRight(90);
}
System.out.println(sp);
System.out.println(p + ", ink: " + sp.inkCount());

Canvas shared = new Canvas(15, 5);
Turtle t1 = new Turtle(shared, 0, 1);
Turtle t2 = new Turtle(t1);
System.out.println(t1.equals(t2));
t1.execute("R 90 F 14");
t2.execute("P o R 135 F 3 L 45 F 11");
System.out.println(t1.equals(t2));
System.out.println(shared);
\end{lstlisting}
\begin{minipage}[t]{0.6\textwidth}
\begin{lstlisting}[style=out]
.............
..***********
..*.........*
..*.*******.*
..*.*.....*.*
..*.*.***.*.*
..*.*.*.*.*.*
..*.*...*.*.*
..*.*****.*.*
..*.......*.*
..*********.*
............*
*************
Turtle at (0,12) heading 0 pen down '*', ink: 79
\end{lstlisting}
\end{minipage}\hfill
\begin{minipage}[t]{0.36\textwidth}
\begin{lstlisting}[style=out]
true
false
...............
o**************
.o.............
..o............
...oooooooooooo
\end{lstlisting}
\end{minipage}

\subsection*{2.4\quad Animation frames (before running)}

The following code records the frames of a small animation in two different
ways. \emph{Before running it}, write down its complete output by hand. Then
run it and compare. Question Q2.4 refers to this code.
\begin{lstlisting}
Canvas f = new Canvas(5, 3);
Turtle ft = new Turtle(f, 0, 1);
ft.turnRight(90);
ArrayList<Canvas> good = new ArrayList<Canvas>();
ArrayList<Canvas> bad = new ArrayList<Canvas>();
for (int i = 0; i < 3; ++i) {
    ft.forward(1);
    good.add(new Canvas(f));
    bad.add(f);
}
for (int i = 0; i < 3; ++i) {
    System.out.println(good.get(i));
    System.out.println(bad.get(i));
    System.out.println();
}
System.out.println(good.get(2).equals(bad.get(0)) + " " + (good.get(2) == bad.get(0))
                   + " " + (bad.get(0) == bad.get(2)));
\end{lstlisting}

\filbreak
\subsection*{2.5\quad Puzzle: the firework}

Starting from \code{new Turtle(new Canvas(11, 11))}, write a program for
\code{execute} with \textbf{at most 25 tokens} that produces exactly this picture
(the centre cell holds an \code{o}):\par
\begin{minipage}{\textwidth}
\begin{lstlisting}[style=out]
*....*....*
.*...*...*.
..*..*..*..
...*.*.*...
...........
****.o.****
...........
...*.*.*...
..*..*..*..
.*...*...*.
*....*....*
\end{lstlisting}
\end{minipage}

\subsection*{2.6\quad Bonus: nested loops}

Extend \code{execute} so that \code{REPEAT} bodies may contain other
\code{REPEAT}s, to any depth. Brackets must then be matched properly. With
\code{new Turtle(new Canvas(19, 5), 1, 3)}, the program
\code{REPEAT 3 [ REPEAT 4 [ F 2 R 90 ] U R 90 F 4 L 90 D ]} must return
\code{true} and draw:
\begin{lstlisting}[style=out]
...................
.***.***.***.......
.*.*.*.*.*.*.......
.***.***.***.......
...................
\end{lstlisting}

\subsection*{2.7\quad Questions}
\begin{questions}
  \item \emph{(Copy semantics.)} The copy constructor of \code{Canvas} copies the
        grid, but the copy constructor of \code{Turtle} shares the canvas.
        Justify \emph{both} choices from what each object represents. Draw the
        memory diagram right after \code{Turtle t2 = new Turtle(t1);} in Test~2.
        What would Test~2 print if \code{Turtle(Turtle)} copied the canvas instead?
  \item \emph{(Identity vs.\ equality.)} \code{Turtle.equals} compares canvases with
        \code{==}, not with \code{Canvas.equals}. Describe a situation where the two
        tests disagree. Argue which one matches the meaning ``these two turtles
        are interchangeable''.
  \item \emph{(Constructors.)} Explain precisely why the naive delegation
        \begin{center}\code{this(canvas, canvas.getWidth() / 2, canvas.getHeight() / 2);}\end{center}
        crashes for a \code{null} canvas. Why can you not simply put an \code{if}
        before \code{this(...)}? How did you solve it? Finally, why does the
        call \code{new Turtle(null)} not even compile, so that one must write
        \code{new Turtle((Canvas) null)}?
  \item \emph{(References and end of life.)} For the code of 2.4: how many
        \code{Canvas} objects exist after the loop? Suppose that after the loop
        we execute \code{f = null; bad = null;}. Which \code{Canvas} objects become
        eligible for garbage collection? Be careful: count every reference
        that still points to each object.
  \item \emph{(Encapsulation.)} Suppose \code{Canvas} offered
        \begin{center}\texttt{public char[][] getCells() \char123\ return cells; \char125}\end{center}
        Write two client
        lines that break a property every other method relies on, and name a
        method that would then crash. Relate this to the
        \code{JeuMorpion.getGrille()} example of Lesson~5.
  \item \emph{(Edge cases.)} What is \texttt{-1 \% 8} in Java? Show step by step how
        your heading update handles \code{turnLeft(450)} from heading $180$, and
        \code{turnRight(-45)} from heading $0$.
\end{questions}

%======================================================================
\exercise{MiniGit: a tiny version-control system}{hard}{5--6 hours}

Version-control systems such as \texttt{git} record \emph{snapshots} of your
files over time. In this exercise you build a miniature one for a single text
document. Every bug in it comes from an object that should have been copied but
was shared, or should have been shared but was copied, so this exercise is a
thorough test of Lessons 8 and 9.

\paragraph{Model.} A \code{Repository} owns a \emph{working copy} (the
\code{Document} that the user edits) and a history of \code{Commit}s. A commit
is a frozen snapshot of the document, with an id, a message and a
\emph{parent} commit (\code{null} for the very first commit). The repository
also knows its \emph{head}, the commit the working copy was last synchronised
with. After a checkout of an old commit, new commits start a new branch, so the
history is a tree. In the test below it looks like this:
\[
  \begin{array}{ccccc}
     & & 2 & \leftarrow & 3 \\
     & \nearrow & & & \\
   1 & & & & \\
     & \nwarrow & & & \\
     & & 4 & \leftarrow & 5
  \end{array}
  \qquad\text{(arrows point from child to parent)}
\]

\paragraph{Warm-up (before coding).} Suppose the \code{Commit} constructor were
written with \code{this.snapshot = content;} instead of making a copy.
Everything else is correct. What would the following code print? Answer
with a memory diagram (this is Q3.1).
\begin{lstlisting}
Repository repo = new Repository();
Document doc = repo.getWorkingCopy();
doc.appendLine("v1");
System.out.println(repo.commit("first"));
doc.setLine(0, "v2");
System.out.println(repo.commit("second"));
System.out.println(repo.log());
\end{lstlisting}

\subsection*{3.1\quad The class \texttt{Document} (mutable)}

A document is a list of lines, stored in a \code{private ArrayList<String>}.
No line is ever \code{null}.

\begin{tabular}{@{}S{0.40\textwidth} D{0.55\textwidth}@{}}
Document() & empty document. \\
Document(String[] lines) & a document with these lines. The array itself must not be kept. \\
Document(Document other) & an independent copy. \\
int lineCount() & number of lines. \\
String getLine(int i) & line \code{i}, or \code{null} if \code{i} is invalid. \\
boolean appendLine(String s) & adds \code{s} at the end. \\
boolean insertLine(int i, String s) & inserts \code{s} before line \code{i}, where $0 \le i \le$ \code{lineCount()}. \\
boolean setLine(int i, String s) & replaces line \code{i}. \\
boolean removeLine(int i) & removes line \code{i}. \\
void replaceContentWith(Document other) & makes \code{this} contain a copy of the lines of \code{other}. Does nothing if \code{other} is \code{null}. Must also work when \code{other == this}. \\
String toString() & each line followed by \nl (empty document: \code{""}). \\
boolean equals(Document o) & same lines in the same order. \\
\end{tabular}

All modifiers return \code{false} and change nothing when the index is invalid
or the line is \code{null}. Otherwise they return \code{true}.

\subsection*{3.2\quad The class \texttt{Commit} (immutable)}

\begin{tabular}{@{}S{0.40\textwidth} D{0.55\textwidth}@{}}
Commit(int id, String message, Document content, Commit parent) & stores an \emph{independent} copy of \code{content}. Later changes to \code{content} must not affect the commit. \\
getId(), getMessage(), getParent() & accessors. \\
Document getSnapshot() & a \emph{new independent copy} of the stored document. \\
boolean hasSameContentAs(Document d) & compares with \code{d} \emph{without} making a copy. \\
int depth() & number of ancestors (the root has depth $0$). \\
boolean isAncestorOf(Commit other) & \code{true} iff this very object is \code{other} or one of its ancestors. \\
String toString() & \code{[3] Serving suggestion (parent 2)} or \code{[1] First version (root)}. \\
\end{tabular}

\subsection*{3.3\quad The class \texttt{Repository}}

A repository has a working document, a head (\code{null} before the first
commit), the list of all its commits, and a counter for the next id. Ids are
$1, 2, 3, \dots$ in creation order. They are assigned by the repository
itself (no \code{static} attribute).

\begin{tabular}{@{}S{0.40\textwidth} D{0.55\textwidth}@{}}
Repository() & empty working document, no commit. \\
Document getWorkingCopy() & the working document \emph{itself} (not a copy), so that clients can edit it. Throughout the life of the repository, it always returns the \emph{same object}. \\
int getHeadId() & id of the head, or $-1$. \\
boolean hasUncommittedChanges() & without a head: whether the working copy is non-empty. Otherwise: whether it differs from the head's snapshot. \\
int commit(String message) & if there is nothing to commit, returns $-1$ and does nothing. Otherwise creates a commit (next id, parent = current head), makes it the head and returns its id. \\
boolean checkout(int id) & returns \code{false} and does nothing if \code{id} is unknown or if there are uncommitted changes. Otherwise the commit becomes the head and the working copy receives its content, \emph{in place}. \\
String log() & the head, then its parent, and so on up to the root, one \code{toString()} per line, separated by \nl, no trailing newline. Without commits: \code{(no commits)}. \\
int mergeBase(int a, int b) & id of the \emph{deepest} commit that is an ancestor of both (every commit is its own ancestor); $-1$ if an id is unknown. \\
String diff(int a, int b) & line diff from snapshot \code{a} to snapshot \code{b} (below); \code{null} if an id is unknown. \\
\end{tabular}

\spec{The diff algorithm}
Let $a_0, \dots, a_{n-1}$ be the lines of the first snapshot and
$b_0, \dots, b_{m-1}$ the lines of the second one. Define $L[i][j]$ as the
length of a longest common subsequence of $a_i, \dots, a_{n-1}$ and
$b_j, \dots, b_{m-1}$:
\[
  L[n][j] = L[i][m] = 0, \qquad
  L[i][j] =
  \begin{cases}
    1 + L[i{+}1][j{+}1] & \text{if } a_i \text{ equals } b_j,\\
    \max\bigl(L[i{+}1][j],\, L[i][j{+}1]\bigr) & \text{otherwise.}
  \end{cases}
\]
Fill the table in a \code{int[n+1][m+1]} array. Then start with $i = j = 0$ and
repeat while $i < n$ or $j < m$:
\begin{enumerate}[itemsep=0pt]
  \item if $i < n$, $j < m$ and $a_i$ equals $b_j$: output \texttt{"\ \ "} $+\,a_i$ (two spaces), then increment $i$ and $j$;
  \item else if $j = m$, or ($i < n$ and $L[i{+}1][j] \geq L[i][j{+}1]$): output \code{"- "} $+\,a_i$, then increment $i$;
  \item else output \code{"+ "} $+\,b_j$, then increment $j$.
\end{enumerate}
Each output line ends with \nl. The tie-break in rule~2
(removals before additions) is part of the specification: it makes the output
unique.

\subsection*{3.4\quad Test}
\begin{lstlisting}
class MiniGitTest {
    public static void main(String[] args) {
        Repository repo = new Repository();
        Document doc = repo.getWorkingCopy();
        System.out.println(repo.log());
        System.out.println(repo.commit("Empty"));

        doc.appendLine("PANCAKES");
        doc.appendLine("200 g flour");
        doc.appendLine("2 eggs");
        doc.appendLine("300 ml milk");
        System.out.println(repo.commit("First version"));
        System.out.println(repo.commit("Nothing new"));

        doc.insertLine(3, "1 pinch of salt");
        doc.setLine(4, "350 ml milk");
        System.out.println(repo.commit("Salt, more milk"));
        doc.appendLine("Serve with maple syrup");
        System.out.println(repo.commit("Serving suggestion"));
        System.out.println(repo.log());

        doc.removeLine(0);
        System.out.println(repo.checkout(1));             // refused
        doc.insertLine(0, "PAN" + "CAKES".substring(0));  // built at run time
        System.out.println(repo.hasUncommittedChanges());
        System.out.println(repo.checkout(1));
        System.out.print(doc);
        System.out.println(doc == repo.getWorkingCopy());

        doc.setLine(0, "VEGAN PANCAKES");
        doc.setLine(2, "1 ripe banana");
        doc.setLine(3, "300 ml oat milk");
        System.out.println(repo.commit("Vegan variant"));
        doc.appendLine("Serve with berries");
        System.out.println(repo.commit("Berries"));
        System.out.println(repo.log());

        System.out.println(repo.mergeBase(3, 5) + " " + repo.mergeBase(2, 3) + " "
                           + repo.mergeBase(4, 4) + " " + repo.mergeBase(1, 9));
        System.out.print(repo.diff(1, 2));
        System.out.println("--");
        System.out.print(repo.diff(3, 5));
        System.out.println("--");
        System.out.println(repo.diff(2, 7));
        System.out.println(repo.checkout(3) + " " + doc.lineCount() + " "
                           + doc.getLine(5) + " " + doc.getLine(6));
    }
}
\end{lstlisting}
Expected output (two columns):

\begin{minipage}[t]{0.45\textwidth}
\begin{lstlisting}[style=out]
(no commits)
-1
1
-1
2
3
[3] Serving suggestion (parent 2)
[2] Salt, more milk (parent 1)
[1] First version (root)
false
false
true
PANCAKES
200 g flour
2 eggs
300 ml milk
true
4
5
[5] Berries (parent 4)
[4] Vegan variant (parent 1)
[1] First version (root)
\end{lstlisting}
\end{minipage}\hfill
\begin{minipage}[t]{0.5\textwidth}
\begin{lstlisting}[style=out]
1 2 4 -1
  PANCAKES
  200 g flour
  2 eggs
- 300 ml milk
+ 1 pinch of salt
+ 350 ml milk
--
- PANCAKES
+ VEGAN PANCAKES
  200 g flour
- 2 eggs
- 1 pinch of salt
- 350 ml milk
- Serve with maple syrup
+ 1 ripe banana
+ 300 ml oat milk
+ Serve with berries
--
null
true 6 Serve with maple syrup null
\end{lstlisting}
\end{minipage}

\subsection*{3.5\quad Play with it: an interactive shell}

Write a class \code{MiniGitShell} whose \code{main} reads commands from the
keyboard (\code{Scanner}) and acts on a \code{Repository}, so that you can version
your own texts:
\begin{center}
\texttt{show, append <text>, insert <i> <text>, set <i> <text>, delete <i>,}\\
\texttt{status, commit <message>, log, checkout <id>, diff <a> <b>, quit}
\end{center}
The exact messages are up to you. The shell must never crash on an invalid
index or id: it reports the failure instead. You may assume that numeric
arguments are well-formed integers.

\subsection*{3.6\quad Questions}
\begin{questions}
  \item \emph{(Warm-up, before coding.)} Answer the warm-up above. Then run it
        on your correct implementation and explain the difference.
  \item \emph{(Interface design.)} \code{getWorkingCopy} returns an internal
        reference, while \code{getSnapshot} returns a copy. Is this
        inconsistent? Justify both choices: which object is \emph{meant} to be
        modified by clients, and which property of the system would a leaking
        \code{getSnapshot} destroy? Write a three-line client snippet that
        demonstrates the damage.
  \item \emph{(In-place update.)} What exactly does the line
        \code{doc == repo.getWorkingCopy()} of the test check?
        Describe the scenario in which a user silently loses edits if
        \code{checkout} did \code{working = c.getSnapshot();} instead.
  \item \emph{(Safe sharing.)} \code{getParent()} returns an internal reference
        without copying it. Is this an encapsulation leak like the one in
        Q3.2? Commit~1 is shared by two branches. Why is this sharing not a
        problem, while sharing a \code{Document} would be?
  \item \emph{(Identity vs.\ equality.)} Why must \code{isAncestorOf} and
        \code{mergeBase} compare commits with \code{==}? Construct two distinct
        commits that an attribute-by-attribute \code{equals} would declare equal,
        and explain why \code{mergeBase} could then return a wrong answer.
  \item \emph{(Strings.)} Why does the test build a line with
        \code{"PAN" + "CAKES".substring(0)}? Which lines of the expected output
        would change if \code{Document.equals} compared lines with \code{==},
        and how?
  \item \emph{(Depth of a copy.)} \code{Document(Document)} copies the
        \code{ArrayList} but shares the \code{String} objects. Why is this deep
        enough? What would have to change if lines were stored as mutable
        objects (for example \code{StringBuilder})?
  \item \emph{(Self-assignment.)} Trace \code{doc.replaceContentWith(doc)} on an
        implementation that first clears its list and then copies the lines of
        \code{other}. What is lost? Explain how your version avoids this.
  \item \emph{(Cost.)} With $c$ commits of a document of about $\ell$ lines,
        how much memory do the snapshots use? What are the time and memory costs
        of \code{diff} on documents of $n$ and $m$ lines? Propose (without
        coding it) a representation of commits that uses less memory, and say
        which public methods would become slower. Explain why the public
        interface of \code{Repository} would not need to change.
\end{questions}

\subsection*{3.7\quad Going further (open-ended, ungraded)}
\begin{itemize}[itemsep=2pt]
  \item \textbf{Named branches.} A \code{Branch} has a name and a head commit
        that moves forward on each commit. Explain why \code{Branch} must be
        mutable while \code{Commit} is immutable, and what that implies for
        sharing and copying.
  \item \textbf{Three-way merge.} Given two heads $A$ and $B$ and their base
        $M = $ \code{mergeBase(A, B)}, combine the changes $M \to A$ and $M \to B$.
        Report a \emph{conflict} when both modify the same line of $M$.
  \item \textbf{Blame.} For each line of the head's snapshot, find the id of the
        commit that introduced it.
\end{itemize}

\end{document}
